% ============================================================
% CAPITOLO 2 — Intelligenza Artificiale nella Pubblica Amministrazione
% JANUS-OP — Tesi di Laurea Magistrale LM-56
% Candidato: Fabrizio Papa | Relatore: Prof. Pasquale Sasso
% Revisione integrata aprile 2026 — aggiornata con letteratura recente
% ============================================================

\chapter{Intelligenza Artificiale nella Pubblica Amministrazione}
\label{chap:ia-pa}

Il capitolo precedente ha delineato il percorso di trasformazione digitale della
Pubblica Amministrazione italiana, evidenziando come il passaggio dal paradigma
dell'\textit{e-Government} a quello del \textit{Digital Government} abbia progressivamente
spostato il baricentro strategico dall'infrastruttura alla capacità decisionale
\parencite{OECDDigitalGovFramework}. L'intelligenza artificiale si colloca al cuore
di questa transizione: non soltanto come strumento tecnologico, ma come fattore
abilitante di una nuova modalità di governo, capace di coniugare efficienza
operativa, trasparenza e \textit{accountability}.

L'analisi che segue adotta una \textit{duplice prospettiva}: \textit{normativa},
per ricostruire il quadro regolatorio europeo e nazionale che disciplina l'uso
dell'IA nel settore pubblico, e \textit{manageriale}, per interpretare le scelte
organizzative attraverso i framework teorici già introdotti nel Capitolo~1
(§\,\ref{chap:digital-government}). In particolare, il \textit{TOE framework}
\parencite{TornatzkyFleischer1990}, le \textit{Dynamic Capabilities}
\parencite{Teece2007} e lo \textit{Strategic Alignment Model}
\parencite{HendersonVenkatraman1993} costituiranno la griglia interpretativa
con cui leggere ciascuna sezione, nel tentativo di superare una lettura puramente
descrittiva dell'adozione dell'IA nella PA a favore di un'analisi critica e
orientata all'azione strategica \parencite{HuergoLora2024}.

Il capitolo è strutturato in cinque sezioni: la prima (§\,\ref{sec:ia-trasformazione})
analizza la trasformazione organizzativa indotta dall'IA nella PA; la seconda
(§\,\ref{sec:ai-act}) esamina il quadro normativo europeo (AI Act, DGA, Data Act);
la terza (§\,\ref{sec:linee-guida-agid}) ricostruisce le linee guida AgID; la quarta
(§\,\ref{sec:modelli-deployment}) analizza i modelli di \textit{deployment} con una
lettura strategica e comparativa; la quinta (§\,\ref{sec:rischi-compliance}) affronta
rischi, compliance e \textit{accountability} algoritmica.

% ============================================================
\section{IA e trasformazione organizzativa nella PA}
\label{sec:ia-trasformazione}

L'intelligenza artificiale non è un semplice strumento di automazione: è una
tecnologia abilitante che ridefinisce le relazioni tra dati, processi decisionali
e strutture organizzative \parencite{HuergoLora2024,ValleCruzMacroMesoMicro2024}.
Nel settore pubblico, questo processo di ridefinizione assume caratteristiche peculiari
legate alla rigidità normativa, alla cultura burocratica, alla molteplicità degli
\textit{stakeholder} e alla natura non-\textit{market} dell'attività amministrativa
\parencite{FernandezRainey2006}.

\textcite{HuergoLora2024} distingue tre gradi di utilizzo della tecnologia digitale
nell'attività amministrativa: l'\textit{e-Government} (comunicazione elettronica con
i cittadini), la \textit{robotizzazione} delle attività amministrative reglamentate e,
al vertice, l'\textit{intelligenza artificiale} propriamente detta, basata su predizioni
algoritmiche e modelli fondazionali. Questa classificazione si articola in continuità
con il percorso evolutivo tracciato nel Capitolo~1 e si rivela cruciale per il progetto
JANUS-OP: il sistema si colloca deliberatamente al terzo livello, orientando l'IA
\textit{on-premise} verso il supporto alle decisioni strategiche della PA piuttosto
che alla mera automatizzazione di procedure standardizzate.

% --- §2.1 ---
\subsection{Framework interpretativi per l'adozione dell'IA nella PA}
\label{subsec:framework-interpretativi}

La letteratura manageriale ha elaborato diversi framework per interpretare l'adozione
di tecnologie innovative nelle organizzazioni. Tre di questi, già introdotti nel
Capitolo~1, risultano particolarmente fecondi per il contesto della PA.

Il \textbf{TOE framework} (\textit{Technology-Organization-Environment}) di
\textcite{TornatzkyFleischer1990} scompone i fattori che influenzano l'adozione
tecnologica in tre dimensioni: il contesto \textit{tecnologico} (caratteristiche e
disponibilità della tecnologia), il contesto \textit{organizzativo} (struttura,
processi e risorse dell'ente) e il contesto \textit{ambientale} (quadro normativo,
mercato, pressioni istituzionali). Applicato all'adozione dell'IA nella PA italiana,
il TOE rivela immediatamente la centralità del contesto ambientale: l'AI Act, il GDPR,
il Perimetro di Sicurezza Cibernetica Nazionale e le direttive AgID --- già analizzati
nel Capitolo~1 nella loro dimensione infrastrutturale --- costituiscono pressioni
ambientali di primo piano che orientano le scelte tecnologiche e organizzative.
\textcite{DelBaroneTOE2025} hanno applicato il TOE framework specificamente alle PA
italiane, confermando che i fattori ambientali risultano determinanti nelle
amministrazioni pubbliche più che nelle imprese private.
\textcite{MadanAshok2022} e \textcite{PulkkinenTOE2025} convergono sulla medesima
conclusione in contesti europei comparati: la pressione regolatoria (AI Act, NIS\,2,
GDPR) orienta le scelte architetturali verso soluzioni che massimizzano il controllo
interno, favorendo il modello \textit{on-premise} per i sistemi ad alto rischio.

Applicato specificamente alla scelta del modello di \textit{deployment on-premise},
il TOE produce la seguente articolazione. Sul piano \textbf{tecnologico}, i fattori
determinanti includono la disponibilità di infrastruttura hardware adeguata (server
GPU, \textit{storage} ad alte prestazioni, rete interna), la compatibilità dei modelli
\textit{open-source} con i sistemi informativi dell'ente, e la maturità della tecnologia
disponibile: la recente proliferazione di modelli \textit{open-weight} di qualità
(Mistral, LLaMA, Falcon) rende tecnicamente praticabile l'\textit{on-premise} anche per
PA di dimensioni medie \parencite{RolandBergerDC2025}. Sul piano \textbf{organizzativo},
i fattori critici comprendono le competenze digitali del personale, la struttura
decisionale dell'ente e la disponibilità di risorse finanziarie per gli investimenti
iniziali. Sul piano \textbf{ambientale}, la normativa (AI Act, GDPR, NIS\,2, Perimetro
Cibernetico) e le linee guida AgID esercitano pressioni coercitive che orientano le PA
verso soluzioni ad alto controllo interno.

Le \textbf{Dynamic Capabilities} di \textcite{Teece2007} --- \textit{sensing}
(percepire opportunità e minacce), \textit{seizing} (coglierle con decisioni rapide)
e \textit{transforming} (riconfigurare le risorse) --- forniscono una chiave di lettura
per interpretare la capacità di un ente pubblico di adottare e sostenere nel tempo
sistemi di IA. Nella PA, le \textit{dynamic capabilities} si manifestano nella capacità
di anticipare i mutamenti del quadro normativo (\textit{sensing}), di selezionare e
implementare soluzioni tecnologiche adeguate (\textit{seizing}) e di riconfigurare i
processi amministrativi (\textit{transforming}) senza incorrere in fenomeni di
\textit{technology lock-in}. Come già osservato nel Capitolo~1, è proprio la
\textit{sensing capability} --- la capacità di leggere i segnali ambientali deboli ---
che distingue le PA che anticipano l'innovazione da quelle che vi reagiscono in ritardo
\parencite{Teece2007}.

Lo \textbf{Strategic Alignment Model} (SAM) di \textcite{HendersonVenkatraman1993},
integrato con gli aggiornamenti di \textcite{CordeiroSAM2019}, completa il quadro:
per essere efficace, la scelta di un sistema di IA \textit{on-premise} deve essere
\textit{allineata} sia con la strategia organizzativa dell'ente (in termini di obiettivi
istituzionali, autonomia decisionale, gestione dei dati sensibili) sia con
l'architettura IT esistente. \textcite{CriadoGilGarcia2019} evidenziano come la mancanza
di allineamento tra strategia IT e strategia organizzativa costituisca la principale
causa di insuccesso nei progetti di digitalizzazione della PA in Europa.

\textcite{MisuracaVanNoort2020}, nella loro indagine JRC su oltre 300 iniziative di IA
nel settore pubblico europeo, rilevano che le amministrazioni con strutture organizzative
flessibili e competenze digitali interne mostrano tassi di successo significativamente
superiori nell'adozione dell'IA --- un dato che conferma empiricamente la centralità
delle \textit{dynamic capabilities} come variabile discriminante.

% --- §2.2 ---
\subsection{Pressioni istituzionali e isomorfismo nell'adozione dell'IA}
\label{subsec:pressioni-istituzionali}

La \textit{Institutional Theory} di DiMaggio e Powell (1983) offre una chiave di
lettura complementare al TOE: le organizzazioni non adottano tecnologie solo per
ragioni di efficienza, ma anche per rispondere a pressioni istituzionali che le spingono
verso forme di \textit{isomorfismo} reciproco. Si distinguono tre meccanismi:

\begin{enumerate}
  \item \textbf{Isomorfismo coercitivo}: deriva da pressioni formali esercitate da
    organizzazioni da cui si dipende. L'AI Act, il GDPR e le direttive AgID esercitano
    pressioni coercitive che inducono uniformità nelle scelte architetturali, orientando
    i sistemi ad alto rischio verso soluzioni \textit{on-premise} dove il controllo è
    pieno e documentabile \parencite{HohmannKollar2025,EnqvistLena2024}.

  \item \textbf{Isomorfismo mimetico}: si manifesta quando le PA, in condizioni di
    incertezza, imitano pratiche di organizzazioni ritenute di successo. Le PA che
    osservano le scelte di amministrazioni pionieristiche (Estonia, Danimarca, AgID per
    l'Italia) tendono a replicarle \parencite{SvardGuerrero2024}.

  \item \textbf{Isomorfismo normativo}: emerge dalla professionalizzazione del campo.
    La formazione degli \textit{AI Officer}, le certificazioni ISO\,42001 e i programmi
    AgID Academy creano standard \textit{de facto} che orientano le scelte indipendentemente
    dall'obbligo formale \parencite{Haynes2026}.
\end{enumerate}

% --- §2.3 ---
\subsection{Dall'automazione al supporto decisionale strategico}
\label{subsec:automazione-decisionale}

Il passaggio dall'automazione al supporto decisionale strategico richiede una
trasformazione profonda delle capacità organizzative che va ben oltre l'implementazione
tecnica. \textcite{CohenLevinthal1990} hanno introdotto il concetto di
\textit{absorptive capacity} --- la capacità di un'organizzazione di riconoscere il
valore di conoscenza esterna, assimilarla e applicarla a fini istituzionali.

Nel contesto di JANUS-OP, l'IA \textit{on-premise} può essere letta come uno strumento
che \textit{potenzia} l'\textit{absorptive capacity} della PA: anziché affidarsi a
\textit{provider} esterni che elaborano e restituiscono \textit{output} opachi, l'ente
sviluppa internamente la capacità di trasformare dati grezzi in conoscenza decisionale.
\textcite{EnqvistLena2024} dimostrano che le amministrazioni con maggiore
\textit{absorptive capacity} interna distinguono con più precisione quando un sistema IA
può essere utilizzato per decisioni autonome (sistemi \textit{rule-based}, basso rischio)
e quando deve limitarsi ad \textit{assistere} il funzionario (sistemi ML ad alto rischio),
replicando esattamente le distinzioni richieste dall'AI Act all'art.~22 GDPR e agli
artt.~14--16 AIA.

\textcite{ShresthaAISDM2024} analizzano come l'IA trasformi il processo di
\textit{strategic decision-making} nella PA, identificando tre modalità di augmentazione
delle capacità decisionali: l'\textit{automation augmentation} (automatizzazione delle
decisioni di routine), la \textit{cognitive augmentation} (potenziamento delle capacità
analitiche del decisore umano) e l'\textit{aggregation augmentation} (sintesi e
ponderazione delle valutazioni di attori multipli, riducendo i \textit{bias} cognitivi
individuali). JANUS-OP opera prevalentemente nelle ultime due modalità, in coerenza con
il principio di \textit{human oversight} dell'AI Act.

% --- §2.4 ---
\subsection{Rischi organizzativi e resistenza al cambiamento}
\label{subsec:rischi-organizzativi}

L'adozione dell'IA nella PA non è esente da resistenze organizzative strutturali.
\textcite{FernandezRainey2006} identificano le condizioni necessarie per il successo del
cambiamento organizzativo nel settore pubblico: supporto della \textit{leadership},
costruzione della coalizione interna, comunicazione della visione, responsabilizzazione
del personale e istituzionalizzazione del cambiamento. Il modello rimanda al framework
di \textcite{Kotter1995}, elaborato in contesti privati ma ampiamente validato nel
settore pubblico.

Le specificità della resistenza al cambiamento nella PA includono: (i) la
\textit{rigidità normativa}, che rende difficile modificare processi codificati in norme
di legge; (ii) la \textit{cultura burocratica}, che privilegia la certezza procedurale
sull'innovazione; (iii) i \textit{vincoli sindacali}, che possono rallentare la
riassegnazione di mansioni automatizzabili; (iv) la \textit{scarsità di competenze
digitali avanzate}, documentata dagli indicatori DESI analizzati nel Capitolo~1 e
confermata da \textcite{MisuracaVanNoort2020} a livello europeo comparato.

\textcite{SvardGuerrero2024}, nella revisione sistematica sull'uso dell'IA nella
gestione dei documenti pubblici in Svezia, Finlandia e Sud Africa, rilevano che la
resistenza organizzativa si manifesta soprattutto nella fase di integrazione tra i nuovi
sistemi IA e i sistemi documentali esistenti (EDMS, protocolli, archivi storici). Le
amministrazioni che hanno superato con successo questa fase sono quelle che hanno
investito nella co-progettazione dei sistemi con il personale operativo.

% --- §2.5 ---
\subsection{IA come leva strategica e decisionale}
\label{subsec:ia-leva-strategica}

La letteratura più recente ha superato la dicotomia tra IA come strumento di efficienza
operativa e IA come rischio organizzativo, convergendo su un terzo paradigma: l'IA come
\textit{leva strategica} che trasforma la natura stessa dell'azione pubblica
\parencite{CriadoGilGarcia2019}.

\textcite{ValleCruzMacroMesoMicro2024} propongono di analizzare l'impatto dell'IA nel
settore pubblico su tre livelli distinti ma interdipendenti. Al livello \textbf{macro},
l'IA ridisegna i regimi di \textit{governance} pubblica, abilitando nuove forme di
\textit{anticipatory governance} --- la capacità di anticipare i problemi collettivi
prima che si manifestino. Al livello \textbf{meso} (ente pubblico), l'IA trasforma il
modello organizzativo, spostando il valore dall'esecuzione procedurale all'analisi e
alla decisione. Al livello \textbf{micro} (processo amministrativo), l'IA accelera la
fase istruttoria, riduce i tempi di risposta al cittadino e aumenta la coerenza delle
decisioni. \textcite{LoreBasile2023} dimostrano empiricamente, attraverso il framework
INTREPID per la PA italiana, come un sistema di IA possa ridurre significativamente i
tempi di verifica della conformità GDPR dei documenti pubblici con un'accuratezza
comparabile a quella degli esperti di dominio.

Il concetto di \textit{ambidexterity organizzativa} di \textcite{OReillyTushman2013}
--- la capacità di essere simultaneamente efficienti nell'operatività corrente
(\textit{exploitation}) e innovativi nell'esplorazione di nuove soluzioni
(\textit{exploration}) --- si rivela particolarmente pertinente per la PA che adotta
l'IA. L'IA \textit{on-premise} serve entrambe le dimensioni: \textit{automatizza} i
processi esistenti liberando risorse umane (\textit{exploitation}) e \textit{abilita}
nuove forme di analisi predittiva e supporto decisionale che prima non erano possibili
(\textit{exploration}). Questa doppia valenza strategica confuta l'obiezione secondo
cui la scelta \textit{on-premise} sarebbe conservativa: al contrario, è la scelta che
massimizza l'\textit{ambidexterity} organizzativa della PA, mantenendo la sovranità sul
dato e aprendo spazi di innovazione sicura \parencite{OReillyTushman2013}.

% ============================================================
\section{Il quadro normativo europeo: AI Act, DGA, Data Act}
\label{sec:ai-act}

Il 2024--2026 ha segnato una svolta epocale nella regolazione dell'IA in Europa: per
la prima volta nella storia, un ordinamento giuridico ha adottato una disciplina organica
dell'intelligenza artificiale come \textit{framework} trasversale che si sovrappone a
tutte le attività regolate dall'IA \parencite{HohmannKollar2025}. L'AI Act si inserisce
in un ecosistema normativo più ampio che comprende il Data Governance Act e il Data Act,
formando la cosiddetta \textit{European Data Strategy} della Commissione.

\subsection{Il Regolamento Europeo sull'Intelligenza Artificiale (AI Act)}
\label{subsec:ai-act-regolamento}

Il Regolamento (UE) 2024/1689 --- l'\textit{AI Act} --- adotta un approccio
\textit{risk-based}: classifica i sistemi di IA in quattro categorie di rischio
(inaccettabile, alto, limitato, minimo) e calibra gli obblighi normativi in funzione
della categoria \parencite{HohmannKollar2025,Haynes2026}.

Per la Pubblica Amministrazione, la classificazione più rilevante è quella dei
\textbf{sistemi ad alto rischio} (Allegato~III), che include espressamente: sistemi di
profiling nella previdenza sociale, sistemi di \textit{law enforcement}, sistemi di
valutazione del rischio fiscale, sistemi per la gestione delle infrastrutture critiche.
Gli obblighi principali per i \textit{deployer} pubblici di sistemi ad alto rischio
includono: (i)~la \textit{Fundamental Rights Impact Assessment} (FRIA), obbligatoria
per le autorità pubbliche (art.~27 AI Act); (ii)~la registrazione nel database UE dei
sistemi ad alto rischio (art.~71); (iii)~la supervisione umana adeguata e documentata
(artt.~14--16); (iv)~la trasparenza verso gli utenti finali (art.~13); (v)~la
conservazione dei \textit{log} e la tracciabilità delle decisioni (art.~12).

\textcite{EnqvistLena2024} analizza la distinzione tra sistemi \textit{rule-based} e
sistemi AI-driven nel contesto della previdenza sociale, dimostrando che i due archetipi
ricadono in categorie normative diverse sia sotto il GDPR (art.~22, decisioni
automatizzate) sia sotto l'AI Act (artt.~14--16 per la supervisione umana). Per la
progettazione di JANUS-OP, questa distinzione è cruciale: il sistema è configurato per
operare come \textit{strumento di supporto alla decisione} ({\textit{human-in-the-loop}})
piuttosto che come motore decisionale autonomo. \textcite{HaitsmaLucas2025} documentano,
attraverso due casi studio olandesi, come la violazione di questo principio abbia
prodotto discriminazioni sistemiche nell'allocazione dei benefici previdenziali.

\textcite{Haynes2026} sottolinea come l'AI Act operi non solo come strumento di
\textit{compliance} ma come \textit{catalizzatore di governance} etica e sostenibile,
operazionalizzando principi come trasparenza, \textit{accountability} e supervisione
umana attraverso un approccio basato sul rischio che si allinea con gli Obiettivi di
Sviluppo Sostenibile dell'ONU (SDG 8, 9, 10, 12 e 16).

\subsection{Data Governance Act e Data Act}
\label{subsec:dga-data-act}

Il \textbf{Data Governance Act} (Regolamento UE 2022/868), in vigore dal settembre
2023, disciplina il riutilizzo dei dati del settore pubblico e istituisce un quadro di
fiducia per la condivisione volontaria dei dati tra soggetti privati. Per la PA, i
profili rilevanti riguardano le condizioni per il riutilizzo di categorie speciali di
dati (sanitari, giudiziari, fiscali) e i requisiti tecnici per l'interoperabilità degli
spazi dati europei --- temi in diretta continuità con il Perimetro di Sicurezza
Cibernetica Nazionale analizzato nel Capitolo~1.

Il \textbf{Data Act} (Regolamento UE 2023/2854), in vigore dal settembre 2025,
introduce il diritto di accesso ai dati generati dall'uso di prodotti connessi e prevede
obblighi di \textit{data portability} per i fornitori di servizi \textit{cloud}. Per i
sistemi di IA della PA, il Data Act ha implicazioni dirette: impone che i dati elaborati
su \textit{cloud} commerciale siano portabili verso infrastrutture alternative ---
argomento normativo a favore della scelta \textit{on-premise}, che elimina alla radice
il problema del \textit{vendor lock-in}. Questo aspetto si ricollega direttamente alla
questione della \textit{tecnological sovereignty} dell'Unione Europea tematizzata nel
Capitolo~1 \parencite{EUParl_TechSovereignty2025}.

\subsection{Interazione tra i tre regolamenti e lettura istituzionale}
\label{subsec:interazione-trittico}

AI Act, DGA e Data Act formano un \textit{ecosistema normativo integrato} che opera su
tre livelli distinti ma complementari \parencite{RohatgiPark2025}. L'AI Act regola i
\textit{sistemi} di IA; il DGA regola la \textit{governance} dei dati; il Data Act
regola i \textit{dati} stessi (diritto di accesso e portabilità).

Riletto attraverso le lenti dell'isomorfismo istituzionale, il trittico regolatorio
rivela la sua natura di campo istituzionale che orienta le scelte delle PA attraverso
tutti e tre i meccanismi: (i)~\textbf{isomorfismo coercitivo}: gli obblighi di FRIA,
registrazione e supervisione umana generano pressione uniforme verso architetture ad alto
controllo --- condizione che il modello \textit{on-premise} soddisfa per costruzione
\parencite{EnqvistLena2024}; (ii)~\textbf{isomorfismo mimetico}: le PA osservano le
scelte architetturali delle PA pionieristiche europee, replicando i modelli di successo
\parencite{SvardGuerrero2024}; (iii)~\textbf{isomorfismo normativo}: la
professionalizzazione del campo (AgID Academy, certificazioni ISO\,42001, figura dell'AI
Officer ex art.~26 AI Act) crea standard \textit{de facto} \parencite{Haynes2026}. Il
risultato è una convergenza nelle pratiche di governance dell'IA nella PA europea verso
soluzioni caratterizzate da alta tracciabilità, supervisione umana documentata e
controllo interno dei dati \parencite{RohatgiPark2025}.

\textcite{HohmannKollar2025} evidenziano che la stratificazione normativa crea zone di
incertezza interpretativa --- in particolare sulla questione se un sistema IA della PA
debba condurre una DPIA (GDPR), una FRIA (AI Act) o entrambe. La risposta corretta,
confermata dalle linee guida EDPB, è che i due strumenti sono complementari e devono
essere condotti separatamente ma in modo coordinato.

% ============================================================
\section{Le linee guida AgID per l'Intelligenza Artificiale}
\label{sec:linee-guida-agid}

L'Agenzia per l'Italia Digitale (AgID) costituisce il principale vettore di isomorfismo
normativo per l'IA nella PA italiana \parencite{AgIDPianoTriennale2024}. La sua funzione
non è solo regolatoria ma anche \textit{abilitante}: attraverso linee guida, standard
tecnici, programmi di formazione e strumenti di valutazione, AgID costruisce il campo
professionale all'interno del quale le PA italiane prendono decisioni sull'IA.

\subsection{Il framework nazionale AgID per l'IA}
\label{subsec:agid-framework}

Le linee guida AgID sull'IA nella PA --- strutturate intorno ai principi di
\textit{affidabilità}, \textit{trasparenza}, \textit{sicurezza},
\textit{privacy by design} e \textit{non discriminazione} --- traducono a livello
operativo i principi dell'AI Act per il contesto nazionale. AgID ha adottato un approccio
\textit{risk-based} coerente con l'AI Act, prevedendo percorsi di conformità differenziati
in funzione del livello di rischio del sistema. Il \textit{Piano Triennale per l'Informatica
nella PA} include specifici obiettivi di adozione dell'IA: sviluppo di servizi predittivi
per la gestione delle risorse, automazione della conformità documentale, supporto alle
decisioni di allocazione delle risorse pubbliche \parencite{AgIDPianoTriennale2024}.

Riletta attraverso il TOE framework, l'azione di AgID opera su tutte e tre le
dimensioni: sul piano \textbf{tecnologico}, fornisce il catalogo delle API e gli
standard di interoperabilità; sul piano \textbf{organizzativo}, produce linee guida,
\textit{toolkit} di valutazione e programmi di formazione (AgID Academy); sul piano
\textbf{ambientale}, raccorda le norme europee con quelle nazionali, riducendo
l'incertezza regolatoria per gli enti locali. Il raccordo con il Piano Triennale
rappresenta inoltre un punto di connessione diretto con la Strategia Cloud Italia
analizzata nel Capitolo~1: il modello \textit{on-premise} proposto da JANUS-OP è
pienamente compatibile con le previsioni di AgID per i dati strategici e ad alto rischio.

\subsection{Applicazioni concrete: conformità GDPR e process mining}
\label{subsec:agid-applicazioni}

Due casi empirici illustrano l'applicazione concreta delle linee guida AgID in ambito IA.

Il primo è il framework \textbf{INTREPID} (\textit{artIficialiNTelligenceforgdpRcompliancEofPublicadmInistrationDocuments})
sviluppato da \textcite{LoreBasile2023} per la PA italiana. INTREPID è un sistema di IA
basato sull'elaborazione del linguaggio naturale (NLP) che verifica automaticamente se i
documenti pubblici in lingua italiana rispettano i requisiti del GDPR. Utilizzando SpaCy
e Tint per il processamento linguistico e un modello di classificazione testuale per
l'identificazione delle violazioni, il framework mostra un'accuratezza comparabile a
quella degli esperti di dominio con tempi di analisi radicalmente inferiori. Crucialmente,
INTREPID opera su infrastruttura interna della PA (\textit{on-premise}), garantendo che
i documenti non vengano trasmessi a server esterni --- requisito implicito del GDPR per
i documenti contenenti dati personali sensibili.

Il secondo caso è lo studio di \textcite{NaiSulis2025} sul \textit{process mining}
applicato agli appalti pubblici europei (dataset TED) e italiani (dati ANAC). Gli autori
dimostrano come le tecniche di \textit{knowledge management} basate su IA consentano di:
(i)~ricostruire automaticamente i processi reali di gara pubblica dai \textit{log} eventi;
(ii)~identificare colli di bottiglia procedurali e deviazioni dai processi normativi
attesi; (iii)~produrre \textit{insight} azionabili per il miglioramento della qualità
procedurale. Entrambi i casi rappresentano l'IA come strumento di supervisione interna
dei processi amministrativi --- precisamente la funzione promossa da AgID.

% ============================================================
\section{Modelli di deployment: on-premise vs cloud vs hybrid}
\label{sec:modelli-deployment}

La scelta del modello di \textit{deployment} per un sistema di IA non è una decisione
tecnica: è una \textit{decisione strategica} che riflette le priorità organizzative, i
vincoli normativi e le \textit{dynamic capabilities} dell'ente
\parencite{HendersonVenkatraman1993}. Tre modelli principali sono disponibili per la PA.

\subsection{Definizioni e tassonomia}
\label{subsec:definizioni-deployment}

Il modello \textbf{\textit{on-premise}} prevede che l'infrastruttura hardware e software
sia interamente sotto il controllo dell'ente. I dati, i modelli IA e i risultati delle
elaborazioni non lasciano mai il perimetro fisico e logico dell'amministrazione. I costi
di implementazione e gestione sono interamente a carico dell'ente, ma il controllo è
massimo e la conformità normativa è garantita per \textit{design}.

Il modello \textbf{\textit{cloud} pubblico} prevede l'utilizzo di risorse computazionali
e servizi IA offerti da \textit{provider} terzi (AWS, Google Cloud, Microsoft Azure,
OVHcloud). Il modello garantisce scalabilità, aggiornamento continuo e riduzione del
\textit{capex}, ma implica la trasmissione di dati a infrastrutture esterne, con
implicazioni rilevanti per la sovranità del dato e la conformità normativa
\parencite{ArunKumar2025}. La questione della sovranità tecnologica, già affrontata nel
Capitolo~1 con riferimento al \textit{CLOUD Act} americano \parencite{CLOUDAct_Wire2025}
e all'iniziativa GAIA-X \parencite{GaiaX2025}, torna qui in tutta la sua rilevanza
operativa per la scelta architetturale.

Il modello \textbf{ibrido} combina elementi dei due precedenti: i dati sensibili e i
sistemi ad alto rischio vengono gestiti \textit{on-premise}, mentre le elaborazioni meno
critiche sono delegate al \textit{cloud}. Questo modello richiede le più alte
\textit{dynamic capabilities}: il \textit{sensing} per identificare quale dato può
transitare sul \textit{cloud}, il \textit{seizing} per configurare le \textit{policy}
di \textit{data governance} dell'architettura ibrida, il \textit{transforming} per
adattarla al mutare del quadro normativo \parencite{Teece2007}.

\subsection{Analisi comparativa per la PA: una lettura TOE}
\label{subsec:analisi-comparativa-deployment}

\begin{table}[h!]
\centering
\caption{Analisi comparativa dei modelli di \textit{deployment} per la PA attraverso il TOE framework}
\label{tab:confronto-deployment}
\renewcommand{\arraystretch}{1.4}
\begin{tabularx}{\textwidth}{lXXX}
\toprule
\textbf{Dimensione} & \textbf{\textit{On-premise}} & \textbf{\textit{Cloud} pubblico} & \textbf{Ibrido} \\
\midrule
Sovranità del dato & Massima & Limitata & Alta (per dati sens.) \\
Conformità GDPR/AI Act & Piena per \textit{design} & Richiede DPA e FRIA & Alta se configurata \\
Costo iniziale (CAPEX) & Elevato & Basso & Medio \\
Scalabilità & Limitata & Massima & Alta \\
\textit{Vendor lock-in} & Nullo & Alto & Medio \\
\textit{Dynamic capabilities} richieste & Medie & Basse & Alte \\
Allineamento NIS\,2/Perimetro & Ottimale & Vincolato & Buono \\
\bottomrule
\end{tabularx}
\end{table}

Riletta attraverso il TOE framework (Tab.~\ref{tab:confronto-deployment}):

\begin{itemize}
  \item Il \textbf{\textit{cloud} pubblico} subisce alta pressione ambientale dalla
    politica \textit{cloud first} del Piano Triennale AgID, ma questa pressione è
    temperata dai vincoli normativi (NIS\,2, Perimetro Cibernetico, GDPR) per i dati
    strategici. Il rischio di disallineamento tecnologico è elevato per i sistemi ad
    alto rischio.
  \item Il modello \textbf{\textit{on-premise}} implica alta complessità organizzativa
    e tecnologica, ma produce il più forte allineamento con il contesto ambientale
    normativo: è la scelta che massimizza la \textit{conformità per design} richiesta
    dall'AI Act \parencite{HohmannKollar2025}.
  \item Il modello \textbf{ibrido} richiede le più alte \textit{dynamic capabilities}
    dell'ente. La gestione della frontiera tra dato \textit{on-premise} e dato in
    \textit{cloud} richiede \textit{sensing} continuo, \textit{seizing} rapido in caso
    di mutamento normativo e \textit{transforming} dell'architettura nel tempo.
\end{itemize}

\subsection{Federated learning come paradigma privacy-preserving}
\label{subsec:federated-learning}

Una soluzione intermedia di crescente rilevanza è il \textit{federated learning} (FL):
un paradigma di apprendimento distribuito in cui i modelli IA vengono addestrati
localmente su ciascun nodo (ente) senza che i dati grezzi vengano mai centralizzati
\parencite{LiLaiWang2023}. Solo i parametri del modello vengono condivisi con un server
di aggregazione, che produce un modello globale migliorato mantenendo la sovranità del
dato in ciascun nodo.

\textcite{LiLaiWang2023} propongono il framework UIFLPP (\textit{Ubiquitous Intelligent
Federated Learning Privacy Protection}) che combina FL con \textit{differential privacy}
e maschere matriciali, raggiungendo un'accuratezza dell'86\% su CIFAR10 con una riduzione
del 5\% dei tempi di addestramento. Per la PA, il FL rappresenta un paradigma
particolarmente interessante per i sistemi inter-amministrativi (es.\ INPS--MEF--Agenzia
delle Entrate): consente di addestrare un modello condiviso su dati fiscali e
previdenziali senza violare il principio di finalità del GDPR. \textcite{ArunKumar2025}
rafforzano ulteriormente questo modello con contratti intelligenti \textit{blockchain}
che garantiscono l'immutabilità dei \textit{log} di trasferimento e la verificabilità
della conformità normativa in ambienti multicloud.

\subsection{Sicurezza informatica e postura digitale dei servizi pubblici}
\label{subsec:sicurezza-deployment}

La scelta del modello di \textit{deployment} non può essere separata dall'analisi della
postura di sicurezza informatica dell'ente --- tema che raccorda direttamente con la
Direttiva NIS\,2 e il Perimetro Cibernetico analizzati nel Capitolo~1.
\textcite{RibeiroFonte2025} hanno valutato la sicurezza di 3.068 piattaforme di servizi
pubblici online in tutti i 193 stati ONU, rilevando che molte continuano a utilizzare
protocolli crittografici obsoleti, certificati mal configurati e \textit{server} non
aggiornati. I paesi con più alta maturità normativa (UE) presentano mediamente una
postura di sicurezza migliore, ma con significativa varianza interna --- a conferma che
la norma non basta: occorre la capacità organizzativa di implementarla
(\textit{dynamic capabilities}).

\textcite{MagnussonIqbal2025}, nella loro revisione sistematica sulla \textit{information
security governance} nel settore pubblico (41 articoli su 1.496 iniziali), identificano
cinque aree critiche: gestione del rischio, \textit{framework} di \textit{accountability},
sicurezza delle reti e-government, piani di mitigazione e allineamento con la strategia
organizzativa. Lo studio rileva che l'aderenza ai framework riconosciuti (ISO/IEC\,27001,
GDPR, NIS\,2) fornisce guida efficace ma che molti enti pubblici difettano
nell'implementazione sostanziale, con scarso utilizzo di KPI di sicurezza e modelli di
maturità. Questa lacuna è esattamente ciò che un sistema \textit{on-premise} come
JANUS-OP può contribuire a colmare, integrando la governance della sicurezza
nell'architettura del sistema fin dalla fase progettuale (\textit{security by design}).

\subsection{Allineamento strategico (SAM) e scelta del modello}
\label{subsec:allineamento-strategico}

Lo Strategic Alignment Model \parencite{HendersonVenkatraman1993,CordeiroSAM2019}
suggerisce che la strategia IT debba essere allineata con la strategia organizzativa
lungo due assi: interno/esterno e strategia/operatività. Applicato alla scelta del
modello di \textit{deployment}:

\begin{enumerate}
  \item Un ente con alta priorità strategica sulla \textit{sovranità del dato} (ministeri,
    agenzie fiscali, forze dell'ordine) dovrebbe scegliere \textit{on-premise} o ibrido
    con forte componente \textit{on-premise} --- posizione di JANUS-OP.
  \item Un ente con alta priorità su scalabilità e \textit{time-to-market} (sportelli
    digitali per il cittadino, servizi informativi) può beneficiare di soluzioni
    \textit{cloud-first}, purché i dati non rientrino nelle categorie sensibili di GDPR
    e AI Act.
  \item Un ente con \textit{portfolio} misto deve adottare un modello ibrido con
    governance chiara delle frontiere tra i due ambienti.
\end{enumerate}

Il progetto JANUS-OP si colloca nel primo scenario: la natura delle decisioni strategiche
supportate (allocazione di risorse, analisi di rischio istituzionale, valutazione di
impatto normativo) richiede il massimo livello di sovranità sul dato. La scelta
\textit{on-premise} non è, in questa prospettiva, conservativa: è la scelta che
massimizza l'allineamento strategico tra architettura IT e missione istituzionale
\parencite{HendersonVenkatraman1993}.

\subsection{Valutazione di opportunità e conformità normativa integrata}
\label{subsec:valutazione-procurement}

La valutazione di un sistema IA per la PA deve integrare tre dimensioni:
opportunità/fattibilità, conformità normativa e piano di monitoraggio continuo. Sul
piano della \textbf{valutazione di opportunità}, gli enti devono considerare il
\textit{Total Cost of Ownership} (TCO) dell'\textit{on-premise} rispetto al modello
\textit{as-a-service}, tenendo conto dei rischi di \textit{lock-in} e della reversibilità
della scelta garantita dal Data Act \parencite{ArunKumar2025}. Sul piano della
\textbf{conformità normativa integrata}, la valutazione deve includere l'\textit{AI
Impact Assessment} previsto dall'AI Act per i sistemi ad alto rischio, la DPIA del GDPR
e la verifica di compatibilità con le linee guida AgID. Sul piano del
\textbf{monitoraggio continuo}, i sistemi IA devono essere soggetti a \textit{audit}
periodici per verificare l'assenza di \textit{model drift} e discriminazione algoritmica
\parencite{HaitsmaLucas2025}.

% ============================================================
\section{Rischi, compliance e accountability algoritmica}
\label{sec:rischi-compliance}

L'adozione dell'IA nella PA introduce rischi di natura tecnica, organizzativa e giuridica
che richiedono un approccio integrato alla \textit{governance} algoritmica. Questa
sezione analizza il quadro normativo applicabile, i rischi specifici della profilazione
algoritmica, gli standard tecnici di riferimento e propone un modello integrato di
\textit{accountability} per la PA.

\subsection{Automazione decisionale: GDPR art.~22 e AI Act}
\label{subsec:automazione-gdpr-aiaact}

Il rapporto tra GDPR e AI Act nella regolazione dell'automazione decisionale nella PA
è uno dei nodi interpretativi più complessi del quadro normativo europeo
\parencite{HohmannKollar2025,EnqvistLena2024}. L'art.~22 GDPR stabilisce il diritto
di ogni interessato a non essere sottoposto a una decisione basata unicamente su
trattamento automatizzato che produca effetti giuridici significativi. Per la PA, questa
disposizione è dirompente: la grande maggioranza delle decisioni amministrative produce
effetti giuridici significativi.

\textcite{EnqvistLena2024} dimostrano che l'art.~22 GDPR e gli artt.~14--16 AI Act
operano in modo diverso per sistemi \textit{rule-based} (trasparenza intrinseca
all'algoritmo) e sistemi ML (la cui trasparenza richiede strumenti di
\textit{explainability} dedicati come SHAP o LIME). Questa distinzione è
cruciale per JANUS-OP: il sistema deve essere configurato per operare come
\textit{human-in-the-loop}, non come motore decisionale autonomo, anche quando la
norma tecnica lo consentirebbe.

\subsection{Discriminazione algoritmica e governance del rischio}
\label{subsec:discriminazione-algoritmica}

La profilazione algoritmica nella PA è uno dei rischi più documentati dell'automazione
decisionale. \textcite{HaitsmaLucas2025} analizzano due agenzie olandesi di previdenza
sociale che hanno adottato sistemi di profilazione algoritmica per contrastare
l'allocazione impropria dei benefici. Lo studio rileva che i sistemi algoritmici hanno
contribuito a istanze di discriminazione sistemica, nonostante la presenza di un quadro
giuridico che richiede espressamente la mitigazione di tali rischi.

Utilizzando i principi di \textit{systemic risk governance} come \textit{framework}
interpretativo, \citeauthor{HaitsmaLucas2025} evidenziano che le agenzie frustrate dalla
complessità sociotecnica tendono a concentrarsi sulla conformità formale (documentazione,
\textit{audit}) senza affrontare la sostanza del rischio discriminatorio. Per la PA
italiana, questo studio costituisce un \textit{warning} specifico: la conformità formale
all'AI Act (FRIA compilata, \textit{log} conservati, supervisore nominato) non è
sufficiente se non è accompagnata da una cultura organizzativa orientata all'equità
algoritmica e da meccanismi di \textit{auditing} sostanziali.

\subsection{Standard tecnici: ISO 42001 e NIST AI RMF}
\label{subsec:standard-tecnici}

I \textit{framework} tecnici per la gestione del rischio dell'IA forniscono alle PA gli
strumenti operativi per tradurre gli obblighi normativi in pratiche di \textit{governance}
concrete. L'\textbf{ISO/IEC\,42001:2023} è il primo standard internazionale per i sistemi
di gestione dell'IA, e definisce i requisiti per l'istituzione, implementazione e
miglioramento continuo di un sistema di gestione dell'IA. Il \textbf{NIST AI Risk
Management Framework} (AI RMF 1.0, 2023) è organizzato in quattro funzioni:
\textit{Govern}, \textit{Map}, \textit{Measure}, \textit{Manage}. I due \textit{framework}
sono complementari e sempre più usati in combinazione anche nel contesto europeo.

\subsection{Information security governance nel settore pubblico}
\label{subsec:security-governance}

\textcite{MagnussonIqbal2025} evidenziano che la conformità ai \textit{framework}
riconosciuti (ISO/IEC\,27001, GDPR, NIS\,2) fornisce orientamento efficace ma che molti
enti pubblici difettano nell'implementazione sostanziale. Questo \textit{gap} tra
conformità formale e sostanziale è particolarmente preoccupante per i sistemi IA, dove
sicurezza informatica e \textit{governance} algoritmica si sovrappongono: una violazione
della sicurezza del sistema IA è contemporaneamente un incidente informatico, una
violazione dei dati personali e un fallimento della \textit{governance} dell'IA.
\textcite{RibeiroFonte2025} confermano questa lettura a scala globale.

\subsection{Un modello integrato di accountability algoritmica per la PA}
\label{subsec:modello-accountability}

Le analisi condotte nelle sezioni precedenti consentono di proporre un modello integrato
di \textit{accountability} algoritmica per la PA che adotta sistemi di IA ad alto rischio.
Il modello si struttura su quattro livelli, derivati dalla sintesi dei \textit{framework}
normativi (AI Act, GDPR) e del modello di IT \textit{Governance} di \textcite{WeillRoss2004},
che distingue le decisioni IT in: principi, architettura, infrastruttura, bisogni
applicativi e investimenti:

\begin{enumerate}
  \item \textbf{Livello~1 --- Governance dei principi}: definizione da parte della
    dirigenza dei valori e criteri che guidano l'uso dell'IA (\textit{AI principles}).
    Corrisponde alla funzione \textit{Govern} del NIST AI RMF e all'art.~26 AI Act.

  \item \textbf{Livello~2 --- Governance dell'architettura}: definizione
    dell'architettura tecnica ammissibile in funzione del rischio del sistema.
    \textcite{WeillRoss2004} identificano questo come il livello più critico nelle
    decisioni di IT \textit{governance}: le scelte architetturali vincolano tutte le
    decisioni successive. Corrisponde alla FRIA dell'art.~27 AI Act.

  \item \textbf{Livello~3 --- Governance dell'infrastruttura}: gestione operativa
    dell'infrastruttura IT (sicurezza, disponibilità, manutenzione). Corrisponde
    alla funzione \textit{Measure} del NIST AI RMF e all'art.~12 AI Act.

  \item \textbf{Livello~4 --- Governance dell'utilizzo}: regolazione dell'uso quotidiano
    del sistema da parte del personale (supervisione, contestazione degli \textit{output},
    formazione continua). \textcite{HaitsmaLucas2025} dimostrano che è proprio questo
    livello il più vulnerabile alle derive discriminatorie. Corrisponde agli artt.~14--16
    AI Act e alla funzione \textit{Manage} del NIST AI RMF.
\end{enumerate}

Il progetto JANUS-OP soddisfa simultaneamente tutti e quattro i livelli: il modello
\textit{on-premise} garantisce la \textit{governance} dell'architettura e
dell'infrastruttura; il \textit{design human-in-the-loop} garantisce la \textit{governance}
dell'utilizzo; l'adozione degli standard ISO\,42001 e NIST AI RMF garantisce la
\textit{governance} dei principi.

% ============================================================
\section*{Sintesi del capitolo}
\addcontentsline{toc}{section}{Sintesi del capitolo}

Il capitolo ha analizzato l'adozione dell'IA nella PA attraverso la duplice prospettiva
normativa e manageriale. I risultati principali possono essere sintetizzati in sei
proposizioni operative, in continuità con le analisi del Capitolo~1:

\begin{enumerate}
  \item L'adozione dell'IA nella PA è un fenomeno istituzionalmente determinato: le
    scelte tecnologiche sono orientate da pressioni coercitive (normativa), mimetiche
    (osservazione dei \textit{best performer}) e normative (professionalizzazione).

  \item Il trittico normativo europeo (AI Act, DGA, Data Act) tende, attraverso
    l'isomorfismo coercitivo, verso architetture caratterizzate da alta sovranità del
    dato e supervisione umana documentata \parencite{HohmannKollar2025,EnqvistLena2024}.

  \item La scelta \textit{on-premise} è strategicamente coerente con la missione
    istituzionale degli enti che trattano dati sensibili ad alto rischio, come
    documentato dall'applicazione del SAM e del TOE framework.

  \item L'IA \textit{on-premise} può aumentare l'\textit{absorptive capacity} della PA,
    trasformando dati grezzi in conoscenza decisionale senza intermediazione esterna
    \parencite{CohenLevinthal1990,LoreBasile2023}.

  \item La conformità formale non è sufficiente: la \textit{governance} algoritmica
    richiede capacità organizzative (\textit{dynamic capabilities}, IT \textit{governance},
    cultura dell'equità) che devono essere sviluppate intenzionalmente
    \parencite{HaitsmaLucas2025,MagnussonIqbal2025}.

  \item L'\textit{ambidexterity} organizzativa \parencite{OReillyTushman2013} è il
    traguardo strategico: l'IA \textit{on-premise} serve simultaneamente
    l'\textit{exploitation} (efficienza operativa) e l'\textit{exploration} (innovazione
    sicura), a condizione che l'ente investa nelle \textit{dynamic capabilities} necessarie.
\end{enumerate}

Queste proposizioni costituiscono il fondamento teorico sul quale il Capitolo~3
costruirà l'analisi della \textit{data governance} e della maturità digitale negli
Atenei, e il Capitolo~4 svilupperà i casi europei e istituzionali di IA \textit{on-premise}
nella PA.

